\documentclass[11pt]{article}
\usepackage[a4paper,margin=25mm]{geometry}
\usepackage{amsmath,amssymb,hyperref}
\title{Distinct dependent Mahler $U_1$-numbers}
\author{ziangni-sys}
\date{Conjecture 00000008326}
\begin{document}
\raggedright
\maketitle
\noindent\textbf{Result.} The universal algebraic-independence assertion is false: two distinct Liouville numbers can differ by one.

\paragraph{Scope.}
Both source versions assert that two $U$-numbers are universally algebraically independent. We refute this conjunct using two numbers in the standard Mahler class $U_1$. No conclusion about the other conjectured laws is needed. In Mahler's classification, the Liouville numbers are precisely the $U_1$-numbers; see Adamczewski--Bugeaud, Definition~1.1 and the following sentence, p.~884~\cite{AB}. This standard classification is used here, independently of the source's separate designation of Liouville numbers as $U_\infty$.

\paragraph{The actual numbers.}
A real number $x$ is Liouville when, for every $N\in\mathbb N$, there are integers $a,b$ with $b>1$ such that
\[
0<\left|x-\frac ab\right|<\frac1{b^N}.
\]
Take the convergent factorial series
\[
L=\sum_{n=0}^{\infty}\frac1{10^{n!}},\qquad M=L+1.
\]
The index starts at zero. This is the base-ten Liouville series with its extra initial $1/10$ term, which does not affect the approximation property. For completeness, its truncation through $n=k$ has denominator $b=10^{k!}$ and integral numerator. For $k\geq\max(1,N)$, its positive tail is bounded by
\[
\sum_{n=k+1}^{\infty}10^{-n!}
\leq \frac{10}{9}\,10^{-(k+1)!}
<10^{-Nk!}=b^{-N}.
\]
Indeed the factorial exponents after $(k+1)!$ increase by at least one, and $((k+1)-N)k!\geq1$. Thus $L$ is Liouville.

For every approximant $a/b$ of $L$, the integer numerator $a+b$ gives
\[
\left|M-\frac{a+b}{b}\right|=
\left|L-\frac ab\right|.
\]
The strict positive error and the denominator are unchanged. Hence $M$ is also Liouville. Both numbers are individually transcendental, and $L\ne M$.

\paragraph{Algebraic dependence.}
The rational polynomial
\[
P(X,Y)=Y-X-1\in\mathbb Q[X,Y]
\]
is nonzero (its value at $(0,0)$ is $-1$), whereas $P(L,M)=0$. The pair is therefore not algebraically independent over $\mathbb Q$. This supplies two distinct $U_1$-numbers, hence two $U$-numbers, contradicting the universal conjunct.

\paragraph{Formal verification.}
The Lean project uses Mathlib's actual \texttt{liouvilleNumber} infinite sum, proves its summability, and applies the verified \texttt{liouville\_liouvilleNumber} approximation theorem. It proves preservation of the full \texttt{Liouville} predicate under translation, individual transcendence, and distinctness. The relation is an actual \texttt{MvPolynomial (Fin 2) \(\mathbb Q\)}; its nonzero value and vanishing evaluation refute Mathlib's \texttt{AlgebraicIndependent} predicate. The concluding existential theorem uses the Liouville approximation characterization of $U_1$, not a formalization of all Mahler classes. No numerical approximations or custom axioms are used.

\begin{thebibliography}{1}
\bibitem{AB} B.~Adamczewski and Y.~Bugeaud, \emph{Transcendence measures for continued fractions involving repetitive or symmetric patterns}, J.~Eur.~Math.~Soc. \textbf{12} (2010), 883--914, p.~884. \href{https://doi.org/10.4171/JEMS/218}{doi:10.4171/JEMS/218}.
\end{thebibliography}
\end{document}
